\documentclass[10pt,a4paper]{amsart}
\usepackage{amsmath,amssymb,amsthm,graphicx}
\usepackage{hyperref}
\numberwithin{equation}{section}

\newtheorem{Theorem}{Theorem}[section]
\newtheorem*{Theorem*}{Theorem}
\newtheorem{Corollary}[Theorem]{Corollary}
\newtheorem{Lemma}[Theorem]{Lemma}
\newtheorem{Proposition}[Theorem]{Proposition}

\theoremstyle{definition}
\newtheorem{Definition}[Theorem]{Definition}
\newtheorem{Note}[Theorem]{Note}
\newtheorem{Example}[Theorem]{Example}
\newtheorem{Remark}[Theorem]{Remark}
\DeclareMathOperator{\Hom}{Hom}
\DeclareMathOperator{\End}{End}
\DeclareMathOperator{\Ker}{Ker}
\DeclareMathOperator{\re}{re}
\DeclareMathOperator{\Ima}{Im}
\DeclareMathOperator{\ev}{ev}
\DeclareMathOperator{\id}{id}
\DeclareMathOperator{\tr}{tr}
\DeclareMathOperator{\ad}{ad}
\DeclareMathOperator{\mult}{mult}
\DeclareMathOperator{\fin}{fin}
\DeclareMathOperator{\str}{str}
\DeclareMathOperator{\row}{row}
\DeclareMathOperator{\col}{col}
\newcommand{\ilim}[1][]{\mathop{\varinjlim}\limits_{#1}}
\newcommand{\plim}[1][]{\mathop{\varprojlim}\limits_{#1}}
\newcommand{\ve}{\varepsilon}
\newcommand{\ts}{\,}
\newcommand{\tss}{\hspace{1pt}}
\newcommand{\on}{\operatorname}
\newcommand{\relmiddle}[1]{\mathrel{}\middle#1\mathrel{}}
\renewcommand{\labelenumii}{\arabic{enumii}).}
\allowdisplaybreaks
\title{Embedding of the tau-fixed central-extension algebra D(2n) into A(2n)}
\author{Mamoru Ueda}
\date{doi:10.3842/SIGMA.2026.089; CC BY 4.0}
\begin{document}
\maketitle
\noindent\textit{Allowed standard prerequisites (editorial).}\begin{itemize}
\item PBW theorem, affine root systems, universal central extensions. Original MRY, GNW and HU bibliography entries; author formulations and uses preserved
\end{itemize}\medskip
\section[Universal central extension of sl(n){[u\^{}\{pm1\},v]}]{Universal central extension of $\boldsymbol{\mathfrak{sl}(n)\bigl[u^{\pm1},v\bigr]}$}
In this section, we recall the work of~\cite{MRY}, which gives a presentation of the universal central extension of $\mathfrak{sl}(n)\bigl[u^{\pm1},v\bigr]$. For a perfect Lie algebra $\mathfrak{g}$, a central extension of $\mathfrak{g}$ is a pair consisting of a~Lie algebra \smash{$\widetilde{\mathfrak{g}}$} and a surjective homomorphism $\pi\colon\widetilde{\mathfrak{g}}\to\mathfrak{g}$ whose kernel is contained in the center of~\smash{$\widetilde{\mathfrak{g}}$}. The central extension $(\widetilde{\mathfrak{g}},\pi)$ is called the universal central extension of $\mathfrak{g}$ if for every central extension $(\mathfrak{e},\psi)$ there exists a unique homomorphism $\phi\colon\widetilde{\mathfrak{g}}\to\mathfrak{e}$ satisfying that $\psi\circ\phi=\pi$.

In~\cite{Kassel}, Kassel gave an explicit presentation of the universal central extension of $\mathfrak{sl}(n)\bigl[u^{\pm1},v\bigr]$.
We set $\bigl(\Omega\bigl(\mathbb{C}\bigl[u^{\pm1},v\bigr]\bigr),d\bigr)$ as the module of differentials of $\mathbb{C}\bigl[u^{\pm1},v\bigr]$ in the same way as in~\cite[Section~2]{MRY}. We take $\{a_i\}$ as a basis of $\mathbb{C}\bigl[u^{\pm1},v\bigr]$. Let $F$ be a free module of $\mathbb{C}\bigl[u^{\pm1},v\bigr]$ whose basis is \smash{$\bigl\{\widetilde{d} a_i\bigr\}$}. We regard $F$ as a 2-sided \smash{$\mathbb{C}\bigl[u^{\pm1},v\bigr]$} module by setting \smash{$a\bigl(\widetilde{d} b\bigr)=\bigl(\widetilde{d} b\bigr)a$} for~any~\smash{$a,b\in\mathbb{C}\bigl[u^{\pm1},v\bigr]$}. We set a linear map \smash{$\widetilde{d} \colon \mathbb{C}\bigl[u^{\pm1},v\bigr]\to F$} as \smash{$\widetilde{d} \bigl(\sum_{i}c_ia_i\bigr)=\sum_{i}c_i\widetilde{d} a_i$} and $K$ as a~submodule of~$F$ generated by \smash{$\widetilde{d}(ab)=(\widetilde{d}a)b+a\bigl(\widetilde{d} b\bigr)$} for $a,b\in\mathbb{C}\bigl[u^{\pm1},v\bigr]$. We denote $F/K$ by~\smash{$\Omega\bigl(\mathbb{C}\bigl[u^{\pm1},v\bigr]\bigr)$}. We also set \smash{$d\colon\mathbb{C}\bigl[u^{\pm1},v\bigr]\to\Omega\bigl(\mathbb{C}\bigl[u^{\pm1},v\bigr]\bigr)$} as a homomorphism induced by \smash{$\widetilde{d}$}.

Let \smash{$\mathfrak{sl}(n)\otimes\mathbb{C}\bigl[u^{\pm1},v\bigr]\oplus\Omega\bigl(\mathbb{C}\bigl[u^{\pm1},v\bigr]\bigr)/d\mathbb{C}\bigl[u^{\pm1},v\bigr]$} be the Lie algebra whose commutator relations are given by
\begin{gather*}
[z_1\otimes a_1,z_2\otimes a_2]=[z_1,z_2]\otimes a_1\cdot a_2+\kappa(z_1,z_2)a_2da_1,\\
\bigl[\mathfrak{sl}(n)\otimes\mathbb{C}\bigl[u^{\pm1},v\bigr],\Omega^1\bigl(\mathbb{C}\bigl[u^{\pm1},v\bigr]\bigr)/d\mathbb{C}\bigl[u^{\pm1},v\bigr]\bigr]=0,\\
\bigl[\Omega^1\bigl(\mathbb{C}\bigl[u^{\pm1},v\bigr]\bigr)/d\mathbb{C}\bigl[u^{\pm1},v\bigr],\Omega^1\bigl(\mathbb{C}\bigl[u^{\pm1},v\bigr]\bigr)/d\mathbb{C}\bigl[u^{\pm1},v\bigr]\bigr]=0.
\end{gather*}
We can naturally define a surjective homomorphism
\[
\pi\colon\ \mathfrak{sl}(n)\otimes\mathbb{C}\bigl[u^{\pm1},v\bigr]\oplus\Omega\bigl(\mathbb{C}\bigl[u^{\pm1},v\bigr]\bigr) / d\mathbb{C}\bigl[u^{\pm1},v\bigr]\to\mathfrak{sl}(n)\otimes\mathbb{C}\bigl[u^{\pm1},v\bigr].
\]
Since the kernel of $\pi$ is $\Omega\bigl(\mathbb{C}\bigl[u^{\pm1},v\bigr]\bigr)/d\mathbb{C}\bigl[u^{\pm1},v\bigr]$, $\pi$ gives a central extension.
Moreover, by~\cite{Kassel}, $\mathfrak{sl}(n)\otimes\mathbb{C}\bigl[u^{\pm1},v\bigr]\oplus\Omega \bigl(\mathbb{C}\bigl[u^{\pm1},v\bigr]\bigr)/d\mathbb{C}\bigl[u^{\pm1},v\bigr]$ together with $\pi$ form the universal central extension of $\mathfrak{sl}(n)\otimes\mathbb{C}\bigl[u^{\pm1},v\bigr]$.

Moody--Rao--Yokonuma~\cite{MRY} gave a new presentation of this universal central extension, which we describe below.
Let us take a matrix \smash{$\bigl(a^A_{i,j}\bigr)_{0\leq i,j\leq n-1}$} as
\begin{equation*}
a^A_{i,j}=\begin{cases}
\hphantom{-}2&\text{if }i=j,\\
-1&\text{if }i=j\pm1,\\
-1&\text{if }(i,j)=(0,n-1),(n-1,0),\\
\hphantom{-}0&\text{otherwise}.
\end{cases}
\end{equation*}
Then, \smash{$\bigl(a^A_{i,j}\bigr)_{0\leq i,j\leq n-1}$} is the Cartan matrix of an affine Lie algebra $\smash{\widehat{\mathfrak{sl}}(n)}=\mathfrak{sl}(n)\otimes\mathbb{C}\bigl[t^{\pm1}\bigr]\oplus\mathbb{C}c$ whose commutator relations are given by
\begin{gather*}
[x\otimes t^u,y\otimes t^s]=[x,y]\otimes t^{u+s}+u\delta_{u+s,0}\tr(xy)c \qquad\text{for }x,y\in\mathfrak{sl}(n),\\
\bigl[c,\widehat{\mathfrak{sl}}(n)\bigr]=0,
\end{gather*}
where $\tr$ is the trace of $\mathfrak{gl}(n)$.
\begin{Definition}%\label{Prop32}
We define $\mathfrak{A}(n)$ as the Lie algebra over $\mathbb{C}$ generated by
\begin{equation*}
\bigl\{x_{i,r}^\pm, h_{i,r}\mid 0\leq i\leq n-1,\,r\in\mathbb{Z}_{\geq0}\bigr\}
\end{equation*}
subject to the following defining relations:
\begin{gather}
[h_{i,r}, h_{j,s}] = 0,\label{Eq2.1}\\
\bigl[x_{i,r}^{+}, x_{j,s}^{-}\bigr] = \delta_{i,j} h_{i, r+s},\label{Eq2.2}\\
\bigl[h_{i,0}, x_{j,r}^{\pm}\bigr] = \pm a^A_{i,j} x_{j,r}^{\pm},\label{Eq2.4}\\
\bigl[h_{i,r+1}, x_{j,s}^{\pm}\bigr] = \bigl[h_{i,r}, x_{j,s+1}^{\pm}\bigr],\label{Eq2.5}\\
\bigl[x_{i, r+1}^{\pm}, x_{j, s}^{\pm}\bigr] = \bigl[x_{i, r}^{\pm}, x_{j, s+1}^{\pm}\bigr],\label{Eq2.8}\\
\sum_{\sigma\in S_{1-a_{i,j}}}\limits\bigl(\ad x_{i,r_{\sigma(1)}}^{\pm}\bigr)\cdots\bigl(\ad x_{i,r_{\sigma(1-a_{i,j})}}^{\pm}\bigr)\bigl(x_{j,s}^{\pm}\bigr)= 0 \qquad\text{if }i \neq j,\label{Eq2.10}
\end{gather}
where $S_l$ is a symmetric group of degree $l$. We also denote the universal enveloping algebra of~$\mathfrak{A}(n)$~by~$A(n)$.
\end{Definition}
By using \eqref{Eq2.8}, we obtain
\begin{equation*}
\bigl[x^\pm_{i,r},x^\pm_{i,s}\bigr]=\frac{1}{2}\bigl(\bigl[x^\pm_{i,r+s},x^\pm_{i,0}\bigr]+\bigl[x^\pm_{i,0},x^\pm_{i,r+s}\bigr]\bigr)=0.%\label{Eq2.11}
\end{equation*}

Moody--Rao--Yokonuma showed that $\mathfrak{A}(n)$ is isomorphic to the universal central extension of~$\mathfrak{sl}(n)\otimes\mathbb{C}\bigl[u^{\pm1},v\bigr]$.
\begin{Theorem}[{\cite[Proposition 3.5]{MRY}}]
There exists an isomorphism
\begin{equation*}
\iota\colon\ \mathfrak{A}(n)\to \mathfrak{sl}(n)\otimes\mathbb{C}\bigl[u^{\pm1},v\bigr]\oplus\Omega\bigl(\mathbb{C}\bigl[u^{\pm1},v\bigr]\bigr)/d\mathbb{C}\bigl[u^{\pm1},v\bigr]
\end{equation*}
given by
\begin{gather*}
x^+_{i,r}=\begin{cases}
E_{i,i+1}v^r&\text{if }i\neq0,\\
E_{n,1}uv^r&\text{if }i=0,
\end{cases}\qquad
x^-_{i,r}=\begin{cases}
E_{i+1,i}v^r&\text{if }i\neq0,\\
E_{1,n}u^{-1}v^r&\text{if }i=0,
\end{cases}
\end{gather*}
where $E_{i,j}$ is a matrix unit of $\mathfrak{gl}(n)$ whose $(p,q)$ component is $\delta_{p,i}\delta_{q,j}$.
\end{Theorem}

Finally, we consider a presentation of the subalgebra of $A(n)$ generated by $\bigl\{x^+_{i,r}\bigr\}$.
We take a~subspace \smash{$\Omega_+\subset\Omega\bigl(\mathbb{C}\bigl[u^{\pm1},v\bigr]\bigr)/d\mathbb{C}\bigl[u^{\pm1},v\bigr]$} spanned by \smash{$\bigl\{u^{r-1}v^sdu\mid r\geq1,\,s\geq0\bigr\}$}.
We set a~subalgebra $A(n)_+\subset A(n)$ generated by \smash{$\bigoplus_{s\geq0}\bigoplus_{\alpha\in\Delta^+}\mathfrak{g}_\alpha v^s$} and $\Omega_+$. Let us take a grading on~$A(n)$ by \smash{${\rm deg}(h_{i,r})={\rm deg}\bigl(x^\pm_{i,r}\bigr)=r+1$} and denote the graded algebra associated with~$A(n)$~by~$\operatorname{gr}A(n)$.
\begin{Theorem}\label{plus}
The associative algebra $A(n)_+$ is isomorphic to the associative algebra generated by $\bigl\{\bar{x}^+_{i,r}\mid 1\leq i\leq n-1,\,r\geq0\bigr\}$ with the relations:
\begin{gather*}
\bigl[\bar{x}^+_{i,r+1},\bar{x}^+_{j,s}\bigr]=\bigl[\bar{x}^+_{i,r},\bar{x}^+_{j,s+1}\bigr],\\
\bigl[\bar{x}^+_{i,r},\bar{x}^+_{j,s}\bigr]=0 \qquad \text{if }a^A_{i,j}=0,\\
\bigl[\bar{x}^+_{i,r},\bigl[\bar{x}^+_{i,s},\bar{x}^+_{j,u}\bigr]\bigr]=0 \qquad \text{if }a^A_{i,j}=-1.
\end{gather*}
\end{Theorem} % RS: pre -> {\rm pre}
Let $\mathfrak{A}(n)_{{\rm pre}}$ be the Lie algebra generated by \smash{$\bigl\{h_{i,r},x^\pm_{i,r}\mid 1\leq i\leq n-1,\,r\geq0\bigr\}$} with the defining relations \eqref{Eq2.1}--\eqref{Eq2.8}. We also define $\mathfrak{A}(n)_{{\rm pre},\pm}$ and $\mathfrak{A}(n)_{{\rm pre},0}$ as Lie subalgebras generated by \smash{$\bigl\{x^\pm_{i,r},h_{i,r}\mid 1\leq i\leq n-1,\,r\geq0\bigr\}$}.
Theorem~\ref{plus} follows from the following lemma.
\begin{Lemma}\qquad
\begin{enumerate}\itemsep=0pt
\item[$(1)$] The associative algebra \raisebox{0.3pt}{$U(\mathfrak{A}(n)_{{\rm pre},\pm})$} is isomorphic to a quotient algebra of the free associative algebra generated by \smash{\raisebox{-0.3pt}{$\bigl\{x^\pm_{i,r}\bigr\}$}} divided by \eqref{Eq2.8}.

\item[$(2)$] Let $I_\pm$ be an ideal of $\mathfrak{A}(n)_{{\rm pre}}$ generated by \eqref{Eq2.10}. Then, $I_-\cap A(n)_{{\rm pre},+}=\varnothing$.
\end{enumerate}
\end{Lemma}
\begin{proof} % RS: не розумію що це, але запис трохи дивно виглядає
\textup{(1)}\ Let $V$ be a quotient algebra of a noncommutative polynomial ring $\mathbb{C}\langle y_{i,d}\mid0\leq i\leq n-1,d\in\mathbb{Z}_{\geq0}\rangle$ divided by the relations $[y_{i,r+1},y_{j,s}]-[y_{i,r},y_{j,s+1}]$. By a direct computation, we obtain a homomorphism
\begin{align*}
\kappa\colon\ \mathfrak{A}(n)_{{\rm pre}}\to\mathfrak{gl}(V)
\end{align*}
by
\begin{align*}
\begin{split}
&\bigl(\kappa\bigl(x^-_{i,d}\bigr)\bigr)(y_{i_1,d_1}\cdots y_{i_r,d_r})=y_{i,d}y_{i_1,d_1}\cdots y_{i_r,d_r},\\
&(\kappa(h_{i,d}))(y_{i_1,d_1}\cdots y_{i_r,d_r})=-\sum_{u=1}^r\limits a^A_{i,i_u}y_{i_1,d_1}\cdots y_{i_{u-1},d_{u-1}}y_{i_u,d+d_u}y_{i_{u+1},d_{u+1}}\cdots y_{i_r,d_r},\\
&\bigl(\kappa\bigl(x^+_{i,d}\bigr)\bigr)(1)=0,\\
&\bigl(\kappa\bigl(x^+_{i,d}\bigr)\bigr)(y_{i_1,d_1}\cdots y_{i_r,d_r})=\delta_{i,i_1}\kappa(h_{i,d+d_1})(y_{i_2,d_2}\cdots y_{i_r,d_r})+y_{i_1,d_1}\bigl(\kappa\bigl(x^+_{i,d}\bigr)\bigr)(y_{i_2,d_2}\cdots y_{i_r,d_r}).
\end{split}
\end{align*}
The homomorphism $\kappa$ induces a surjective homomorphism $\Xi\colon \mathfrak{A}(n)_{{\rm pre},-}\to V$ by $\Xi(x)=(\kappa(x))1$.
By a direct computation, we can also construct a homomorphism
\begin{align*}
\xi\colon\ V\to \mathfrak{A}(n)_{{\rm pre},-},\qquad y_{i,d}\mapsto x^-_{i,d}.
\end{align*}
Since $\xi\circ\Xi$ is an identity map, $\Xi$ is bijective. Then, we find that $U(\mathfrak{A}(n)_{{\rm pre},-})$ is isomorphic to a quotient algebra of the free associative algebra generated by $\{x^-_{i,r}\}$ divided by \eqref{Eq2.8}. Since~$\mathfrak{A}(n)_{{\rm pre},+}$ is isomorphic to $\mathfrak{A}(n)_{{\rm pre},-}$, $U(\mathfrak{A}(n)_{{\rm pre},+})$ is also isomorphic to a quotient algebra of the free associative algebra generated by $\bigl\{x^+_{i,r}\mid r\geq0\bigr\}$ divided by \eqref{Eq2.8}.

\textup{(2)}\ It is enough to show that $I_-$ is contained in $\mathfrak{A}(n)_{{\rm pre},-}$. This follows from the fact that $\ad\bigl(x^+_{i,r}\bigr)$ acts on the left-hand side of \eqref{Eq2.10} trivially.
\end{proof}


\section[A central extension of sl(2n){[u\^{}\{pm1\},v]}\^{}tau]{A central extension of $\boldsymbol{\mathfrak{sl}(2n)\bigl[u^{\pm1},v\bigr]^\tau}$}\label{sec:5}
We give another involution of $\mathfrak{sl}(2n)$. Set $I_n=\{\pm1,\pm2,\dots,\pm n\}$ and take $\{E_{i,j}\}_{i,j\in I_n}$ as the set of matrix units in $\mathfrak{gl}(2n)$. We also define $\tau$ as an automorphism of $\mathfrak{gl}(2n)\bigl[u^{\pm1},v\bigr]$ by $\tau(E_{i,j}u^rv^s)=-(-1)^sE_{-j,-i}u^rv^s$. By definition of $\tau$, $\mathfrak{sl}(2n)\bigl[u^{\pm1},v\bigr]$ is preserved by $\tau$. In the next section, we give a central extension of $\mathfrak{sl}(2n)\bigl[u^{\pm1},v\bigr]^\tau$.

For $r\geq0$, we define matrices \smash{$\bigl(a^{2r}_{i,j}\bigr)_{0\leq i\leq n,0\leq j\leq n}$} and \smash{$\bigl(a^{2r+1}_{i,j}\bigr)_{1\leq i\leq n-1,0\leq j\leq n}$} as follows:
\begin{align*}
a^{2r}_{i,j}&=\begin{cases}
2\delta_{i,j}-\delta_{i+1,j}-\delta_{i,j+1}&\text{if }3\leq i\leq n-3,\\
2\delta_{j,n-2}-\delta_{j,n-3}-\delta_{j,n-1}-\delta_{j,n}&\text{if }i=n-2,\\
2\delta_{j,n-1}-\delta_{j,n-2}&\text{if }i=n-1,\\
2\delta_{j,n}-\delta_{j,n-2}&\text{if }i=n,\\
2\delta_{j,2}-\delta_{j,0}-\delta_{j,1}-\delta_{j,3}&\text{if }i=2,\\
2\delta_{j,1}-\delta_{j,2}&\text{if }i=1,\\
2\delta_{j,0}-\delta_{j,2}&\text{if }i=0,
\end{cases}\\
a^{2r+1}_{i,j}&=\begin{cases}
a^{2r}_{i,j}&\text{if }2\leq i\leq n-2,\\
2\delta_{j,n-1}+2\delta_{j,n}-\delta_{j,n-2}&\text{if }i=n-1,\\
2\delta_{j,1}+2\delta_{j,0}-\delta_{j,2}&\text{if }i=1.
\end{cases}
\end{align*}
We note that \smash{$\bigl(a^{2r}_{i,j}\bigr)_{i,j=1}^n$} is the Cartan matrix of $\widehat{\mathfrak{so}}(2n)$.
\begin{Definition}
We define $\mathfrak{D}(2n)$ as the Lie algebra generated by
\begin{equation*}
\bigl\{H_{i,r},X^\pm_{i,r},H_{n,2r},H_{0,2r},X^\pm_{-1,2r+1},X^\pm_{n+1,2r+1} \mid 1\leq i\leq n-1,\, r\in\mathbb{Z}_{\geq0}\bigr\}
\end{equation*}
with the following commutator relations:
\begin{gather}
[H_{i,r},H_{j,s}]=0, \label{5111}\\
\bigl[H_{i,r},X^\pm_{j,s}\bigr]=\pm a^r_{i,j}X^\pm_{j,r+s}\qquad \text{for }0\leq j\leq n,\label{5112}\\
\bigl[X^+_{i,r},X^-_{j,s}\bigr]=\begin{cases}
\delta_{i,j}H_{i,r+s}& \text{if} \ 0\leq i,j\leq n,\\
& \hphantom{\text{for}}\! (i,j)\neq(n-1,n),(n,n-1),(0,1),(1,0),(n,n),(0,0),\\
H_{n,r+s}& \text{if }i=j=n\text{ and }r+s\text{ is }even,\\
H_{n-1,r+s}& \text{if }i=j=n\text{ and }r+s\text{ is }odd,\\
H_{0,r+s}& \text{if }i=j=0\text{ and }r+s\text{ is }even,\\
H_{1,r+s}& \text{if }i=j=0\text{ and }\text{ and }r+s\text{ is }odd,
\end{cases}\!\!\! \label{5113}\\
\bigl[X^+_{n-1,r},X^-_{n,s}\bigr]=-\delta_{r+s,{\rm odd}}X^-_{n+1,r+s},\qquad \bigl[X^+_{n,r},X^-_{n-1,s}\bigr]=-\delta_{r+s,{\rm odd}}X^+_{n+1,r+s},\label{5114}\\
\bigl[X^+_{0,r},X^-_{1,s}\bigr]=-\delta_{r+s,{\rm odd}}X^+_{-1,r+s},\qquad \bigl[X^+_{1,r},X^-_{0,s}\bigr]=-\delta_{r+s,{\rm odd}}X^-_{-1,r+s}, \label{5115}\\
\bigl[X^\pm_{i,r},X^+_{n+1,s}\bigr]=0,\qquad\bigl[X^\pm_{i,r},X^-_{n+1,s}\bigr]=0 \qquad\text{if }1\leq i\leq n-1, \label{5116}\\
\bigl[X^\pm_{i,r},X^+_{-1,s}\bigr]=0,\qquad\bigl[X^\pm_{i,r},X^-_{-1,s}\bigr]=0 \qquad\text{if }2\leq i\leq n+1,\label{5117}\\
\bigl[X^+_{n-1,r},X^+_{n+1,s}\bigr]=2X^+_{n,r+s},\qquad \bigl[X^-_{n-1,r},X^+_{n+1,s}\bigr]=0,\label{5118}\\
\bigl[X^+_{n-1,r},X^-_{n+1,s}\bigr]=0,\qquad \bigl[X^-_{n-1,r},X^-_{n+1,s}\bigr]=-2X^-_{n,r+s},\label{5119}\\
\bigl[X^+_{n,r},X^+_{n+1,s}\bigr]=0,\qquad \bigl[X^-_{n,r},X^+_{n+1,s}\bigr]=-2X^{r+s}_{n-1,r+s}, \label{5120}\\
\bigl[X^+_{n,r},X^-_{n+1,s}\bigr]=2X^+_{n-1,r+s},\qquad \bigl[X^-_{n,r},X^-_{n+1,s}\bigr]=0, \label{5121}\\
\bigl[X^+_{0,r},X^+_{-1,s}\bigr]=0,\qquad \bigl[X^-_{0,r},X^+_{-1,s}\bigr]=-2X^-_{1,r+s}, \label{5122}\\
\bigl[X^+_{0,r},X^-_{-1,s}\bigr]=2X^+_{1,r+s},\qquad \bigl[X^-_{0,r},X^-_{-1,s}\bigr]=0, \label{5123}\\
\bigl[X^+_{1,r},X^+_{-1,s}\bigr]=2X^+_{0,r+s},\qquad \bigl[X^-_{1,r},X^+_{-1,s}\bigr]=0, \label{5124}\\
\bigl[X^+_{1,r},X^-_{-1,s}\bigr]=0,\qquad \bigl[X^-_{1,r},X^-_{-1,s}\bigr]=-2X^-_{0,r+s}, \label{5125}\\
\bigl[X^\pm_{i,r},X^\pm_{i,s}\bigr]=0, \label{5126}\\
\bigl[X^\pm_{i,r+1},X^\pm_{j,s}\bigr]=\bigl[X^\pm_{i,r},X^\pm_{j,s+1}\bigr]\quad \text{if }0\leq i\leq j\leq n,(i,j)\neq(0,1),(n-1,n),\label{5127}\\
\bigl[X^\pm_{0,r+1},X^\pm_{1,s}\bigr]=-\bigl[X^\pm_{0,r},X^\pm_{1,s+1}\bigr],\qquad \bigl[X^\pm_{n-1,r+1},X^\pm_{n,s}\bigr]=-\bigl[X^\pm_{n-1,r},X^\pm_{n,s+1}\bigr], \label{5128}\\
\bigl[X^\pm_{i,r},X^\pm_{j,s}\bigr]=0\quad \text{if }a^0_{i,j}=0,0\leq i\neq j\leq n,(i,j)\neq(0,1),(1,0),(n-1,n),(n,n-1),\!\!\!\label{5129}\\
\bigl[X^\pm_{i,r},\bigl[X^\pm_{i,s},X^\pm_{j,u}\bigr]\bigr]=0\qquad \text{if }a^0_{i,j}=-1,\label{5130}\\
\bigl[X^\pm_{n-1,0},\bigl[X^\pm_{n-1,0},X^\pm_{n,2r+1}\bigr]\bigr]=0,\label{5131}\\
\bigl[X^\pm_{n,0},\bigl[X^\pm_{n,0},X^\pm_{n-1,2r+1}\bigr]\bigr]=0,\label{5132}\\
\bigl[X^\pm_{1,0},\bigl[X^\pm_{1,0},X^\pm_{0,2r+1}\bigr]\bigr]=0,\label{5133}\\
\bigl[X^\pm_{0,0},\bigl[X^\pm_{0,0},X^\pm_{1,2r+1}\bigr]\bigr]=0,\label{5134}\\
\bigl[X^+_{1,1},\bigl[X^-_{1,0},(E_{1,1}-E_{2,2})t^u\bigr]\bigr]-\bigl[X^+_{1,0}, \bigl[X^-_{1,1},(E_{1,1}-E_{2,2})t^u\bigr]\bigr]=0. \label{5135}
\end{gather}
We also denote the universal enveloping algebra of $\mathfrak{D}(2n)$ by $D(2n)$.
\end{Definition}
Set \smash{$f^{r,s}_{i,j}=(E_{i,j}-(-1)^sE_{-j,-i})u^rv^s\in\mathfrak{sl}(2n)\bigl[u^{\pm1},v\bigr]$}. Then, \smash{$\mathfrak{sl}(2n)\bigl[u^{\pm1},v\bigr]^\tau$} is spanned by \smash{$\bigl\{f^{r,s}_{i,j}\mid i,j\in I_n,\,r\in\mathbb{Z},\,s\in\mathbb{Z}_{\geq0}\bigr\}$}. By a direct computation, we can construct a homomorphism \smash{$\pi\colon \mathfrak{D}(2n)\to\mathfrak{sl}(2n)\bigl[u^{\pm1},v\bigr]^\tau$} determined by
\begin{gather*}
H_{i,r}=\begin{cases}
f^r_{i,i}v^r-f^r_{i+1,i+1}v^r&\text{if }1\leq i\leq n-1,\\
f^r_{n-1,n-1}v^r+f^r_{n,n}v^r&\text{if }i=n\text{ and $r$ is even},\\
-\bigl(f^r_{1,1}v^r+f^r_{2,2}v^r\bigr)&\text{if }i=0\text{ and $r$ is even},
\end{cases}\\
X^+_{i,r}=\begin{cases}
f^r_{i,i+1}v^r&\text{if }1\leq i\leq n-1,\\
f^r_{n-1,-n}v^r&\text{if }i=n,\\
(-1)^rf^r_{-2,1}uv^r&\text{if }i=0,\\
f^r_{n,-n}v^r&\text{if }i=n+1\text{ and $r$ is odd},\\
f^r_{-1,1}uv^r&\text{if }i=-1\text{ and $r$ is odd},
\end{cases}\\
X^-_{i,r}=\begin{cases}
f^r_{i+1,i}v^r&\text{if }1\leq i\leq n-1,\\
f^r_{-n,n-1}v^r&\text{if }i=n,\\
(-1)^rf^r_{1,-2}u^{-1}v^r&\text{if }i=0,\\
f^r_{-n,n}v^r&\text{if }i=n+1\text{ and $r$ is odd},\\
f^r_{1,-1}u^{-1}v^r&\text{if }i=-1\text{ and $r$ is odd}.
\end{cases}
\end{gather*}
Next, we show that $\mathfrak{D}(2n)$ is a central extension of $\mathfrak{sl}(2n)\bigl[u^{\pm1},v\bigr]^\tau$.

Let $\{\alpha_i\mid0\leq i\leq n\}$ be the set of simple roots of $\widehat{\mathfrak{so}}(2n)$, let $W$ be the Weyl group of $\widehat{\mathfrak{so}}(2n)$, and let~$\Delta$ be the set of roots of $\widehat{\mathfrak{so}}(2n)$. We also denote
\begin{gather*}
Q=\sum_{i=0}^n \mathbb{Z}\alpha_i,Q_+=\sum_{i=0}^n \mathbb{Z}_{\geq0}\alpha_i\setminus\{0\},\qquad Q_-=-Q_+\qquad \text{and}\\ \delta=\alpha_0+\alpha_1+\alpha_{n-1}+\alpha_n+\sum_{j=2}^{n-2} \alpha_j.
\end{gather*} We also take a subset of $Q$ as follows:
\begin{align*}
\Delta^{{\rm ex}}=\Delta\cup\Biggl\{&\pm\Bigg(2\sum_{j=i}^{n-2}\alpha_j+\alpha_{n-1}+\alpha_n\Biggr)+s\delta,\\
&\pm(\alpha_{n-1}-\alpha_n)+s\delta,\pm(\alpha_0-\alpha_1)+s\delta\mid 1\leq i\leq n-2,\,s\in\mathbb{Z}\Biggr\}.
\end{align*}
We define a degree of $\mathfrak{D}(2n)$ by $Q\times\mathbb{Z}_{\geq0}$ by assigning generators as follows:
\begin{gather*}
{\rm deg}(H_{i,r})=(0,r),\qquad {\rm deg}(X^\pm_{i,r})=\begin{cases}
(\pm\alpha_i,r)&\text{if }0\leq i\leq n,\\
(\pm(\alpha_{n}-\alpha_{n-1}),r)&\text{if }i=n+1,\\
(\pm(\alpha_{0}-\alpha_1),r)&\text{if }i=0.
\end{cases}
\end{gather*}
We denote the space of elements of degree $(\alpha,k)$ in $\mathfrak{D}(2n)$ by \smash{$\mathfrak{D}(2n)_\alpha^k$}. The following lemma can be proved in the same way as~\cite[Lemma 3.1]{MRY}.
\begin{Lemma}\label{Lem128}\qquad
\begin{enumerate}\itemsep=0pt
\item[$(1)$] Let $\mathfrak{s}_\pm$ be the subspace of $\mathfrak{D}(2n)$ spanned by \smash{$\bigl\{\prod_{i=1}^v\ad\bigl(X^\pm_{u_i,r_i}\bigr)X^\pm_{u_{v+1},r_{v+1}}\bigr\}$}
and~$\mathfrak{h}$ as the space spanned by $\{H_{i,r}\}$. Then, $\mathfrak{D}(2n)$ can be decomposed to
\begin{align*}
\mathfrak{s}^+& \oplus\mathfrak{h}\oplus\mathfrak{s}^-\oplus \bigoplus_{r\in\mathbb{Z}_{\geq0}}\mathbb{C}X^+_{n+1,2r+1}\oplus \bigoplus_{r\in\mathbb{Z}_{\geq0}}\mathbb{C}X^-_{n+1,2r+1}\\
& \oplus \bigoplus_{r\in\mathbb{Z}_{\geq0}}\mathbb{C}X^+_{-1,2r+1}\oplus \bigoplus_{r\in\mathbb{Z}_{\geq0}}\mathbb{C}X^-_{-1,2r+1}.
\end{align*}
\item[$(2)$] If $\alpha\in Q_{\pm}$, $\mathfrak{D}(2n)_\alpha^k$ is spanned by
\begin{equation*}
\Biggl\{\prod_{i=1}^v\ad\bigl(X^\pm_{u_i,r_i}\bigr)X^\pm_{u_{v+1},r_{v+1}}\:\bigg|\: \sum_{i=1}^{v+1}\alpha_{u_i}=\alpha,\,\sum_{i=1}^{v+1}r_i=k\Biggr\}.
\end{equation*}
\end{enumerate}
\end{Lemma}
For $0\leq i\leq n$, we set
\[ s_i\in W \qquad \text{ as } s_i(\alpha_j)=\begin{cases}
-\alpha_i&\text{if }i=j,\\
\alpha_i+\alpha_j&\text{if }a^0_{i,j}=-1,\\
\alpha_j&\text{otherwise.}
\end{cases}
\]
\begin{Lemma}%\label{root}
By the action of the Weyl group, $\Delta^{{\rm ex}}$ is generated by $\{\alpha_j\mid0\leq j\leq n\}$ and $\pm(\alpha_{n-1}-\alpha_n)$, $\pm(\alpha_0-\alpha_1)$.
\end{Lemma}
\begin{proof}
Any root of $\Delta$ can be derived from $\{\alpha_j\mid0\leq j\leq n\}$ by the Weyl group. Thus, it is enough to show that \smash{$\pm\bigl(2\sum_{j=i}^{n-2}\alpha_j+\alpha_{n-1}+\alpha_n\bigr)+s\delta$} and $\pm(\alpha_{n-1}-\alpha_n)+s\delta$ can be derived from~$\pm(\alpha_{n-1}-\alpha_n)$,~$\pm(\alpha_0-\alpha_1)$.

First, we consider the case where $s$ is even.
By the definition of $s_i$, we obtain the following relations:
\begin{align*}
&s_n(\alpha_{n-1}-\alpha_n+s\delta)=\alpha_{n-1}+\alpha_n+s\delta,\\
&s_{n-1}(\alpha_{n-1}+\alpha_n+s\delta)=-\alpha_{n-1}+\alpha_n+s\delta,\\
&s_{i-1}\Biggl(2\sum_{j=i}^{n-2}\alpha_j+\alpha_{n-1}+\alpha_n+s\delta\Biggr)=2\sum_{j=i-1}^{n-2}\alpha_j+\alpha_{n-1}+\alpha_n+s\delta,\\
&s_{0}\Biggl(2\sum_{j=1}^{n-2}\alpha_j+\alpha_{n-1}+\alpha_n+s\delta\Biggr)=-\alpha_{0}-\alpha_1+(s+1)\delta,\\
&s_0(-\alpha_{0}-\alpha_1+(s+1)\delta)=\alpha_{0}-\alpha_1+(s+1)\delta,\\
&s_1(\alpha_0-\alpha_1+s\delta)=\alpha_{0}+\alpha_1+s\delta=-2\sum_{j=1}^{n-2}\alpha_j-\alpha_{n-1}-\alpha_n+(s+1)\delta,\\
&s_{i}\Biggl(-2\sum_{j=i}^{n-2}\alpha_j-\alpha_{n-1}-\alpha_n+(s+1)\delta\Biggr)=-2\sum_{j=i+1}^{n-2}\alpha_j-\alpha_{n-1}-\alpha_n+(s+1)\delta,\\
&s_{n-1}(-\alpha_{n-1}-\alpha_n+(s+1)\delta)=\alpha_{n-1}-\alpha_n+(s+1)\delta.
\end{align*}
Then, by these relations, \smash{$\bigl(2\sum_{j=i}^{n-2}\alpha_j+\alpha_{n-1}+\alpha_n\bigr)+s\delta$}, \smash{$-\bigl(2\sum_{j=i}^{n-2}\alpha_j+\alpha_{n-1}+\alpha_n\bigr)+(s+1)\delta$}, $\pm(\alpha_{n-1}-\alpha_n)+s\delta$ and $\pm(\alpha_0-\alpha_1)+s\delta$ can be derived from $\alpha_{n-1}-\alpha_n$ and $\alpha_0-\alpha_1$ for $s\geq0$.
Similarly to the positive case, $\bigl(2\sum_{j=i}^{n-2}\alpha_j+\alpha_{n-1}+\alpha_n\bigr)-(s+1)\delta$, $-\bigl(2\sum_{j=i}^{n-2}\alpha_j+\alpha_{n-1}+\alpha_n\bigr)-s\delta$, $\pm(\alpha_{n-1}-\alpha_n)-(s+1)\delta$ and $\pm(\alpha_0-\alpha_1)-(s+1)\delta$ can be derived from $-(\alpha_{n-1}-\alpha_n)$ and $-(\alpha_0-\alpha_1)$ for~${s\geq0}$.
\end{proof}
By the relations \eqref{5129}--\eqref{5134}, we can define an automorphism of $D(2n)$ by
\begin{gather*}
\tau_i=\exp\bigl(\ad\bigl(X^+_{i,0}\bigr)\bigr)\exp\bigl(-\ad\bigl(X^-_{i,0}\bigr)\bigr)\exp\bigl(\ad\bigl(X^+_{i,0}\bigr)\bigr).
\end{gather*}
By a direct computation, we obtain the following lemma.
\begin{Lemma}\label{Weyl}
The following relations hold for $1\leq i\leq n-2$:
\begin{gather*}
\tau_i(X^+_{j,r})=\begin{cases}
-X^-_{i,r}&\text{if }i=j,\\
\bigl[X^+_{i,0},X^+_{j,r}\bigr]&\text{if }a^0_{i,j}=-1,\\
(-1)^rX^+_{j,r}&\text{if }(i,j)=(0,1),(1,0),(n-1,n),(n,n-1),\\
X^+_{i,r}&\text{otherwise},
\end{cases}\\
\tau_i(X^-_{j,r})=\begin{cases}
-X^+_{i,r}&\text{if }i=j,\\
-\bigl[X^-_{i,0},X^-_{j,r}\bigr]&\text{if }a^0_{i,j}=-1,\\
(-1)^rX^-_{j,r}&\text{if }(i,j)\neq(0,1),(1,0),(n-1,n),(n,n-1),\\
X^+_{j,r}&\text{otherwise}.
\end{cases}
\end{gather*}
\end{Lemma}
\begin{proof}
We only consider the case where $(i,j)=(n-1,n)$ when $r$ is odd. The other cases can be proved in a similar way. First, we prove the $+$ case. By \eqref{5131}, we have
\begin{align*}
&\exp\bigl(\ad\bigl(X^+_{n-1,0}\bigr)\bigr)X^+_{n,2r+1}=X^+_{n,2r+1}+\bigl[X^+_{n-1,0},X^+_{n,2r+1}\bigr]=X^+_{n,2r+1}+\bigl[X^+_{n-1,0},X^+_{n,2r+1}\bigr].
\end{align*}
By \eqref{5114}, \eqref{5112} and \eqref{5118}, we obtain
\begin{align*}
&-\ad\bigl(X^-_{n-1,0}\bigr)\bigl(X^+_{n,2r+1}+\bigl[X^+_{n-1,0},X^+_{n,2r+1}\bigr]\bigr)\\
&\qquad{}=-\bigl[X^-_{n-1,0},X^+_{n,2r+1}\bigr]-\bigl[X^-_{n-1,0},\bigl[X^+_{n-1,0},X^+_{n,2r+1}\bigr]\bigr]\\
&\qquad{}=-X^+_{n+1,2r+1}+\bigl[h_{n-1,0},X^+_{n,2r+1}\bigr]-\bigl[X^+_{n-1,0},\bigl[X^-_{n-1,0},X^+_{n,2r+1}\bigr]\bigr]\\
&\qquad{}=-X^+_{n+1,2r+1}-\bigl[X^+_{n-1,0},X^+_{n+1,2r+1}\bigr]=X^+_{n+1,2r+1}-2X^+_{n,2r+1}.
\end{align*}
By \eqref{5118} and \eqref{5114}, we have
\begin{align*}
-\frac{1}{2}\ad\bigl(X^-_{n-1,0}\bigr)\bigl(-X^+_{n+1,2r+1}-2X^+_{n,2r+1}\bigr)
&=0+X^+_{n+1,2r+1}.
\end{align*}
Thus, we obtain
\begin{align*}
\exp\bigl(-\ad\bigl(X^-_{n-1,0}\bigr)\bigr)\exp\bigl(\ad\bigl(X^+_{n-1,0}\bigr)\bigr)X^+_{n,2r+1}&=-X^+_{n,2r+1}+\bigl[X^+_{n-1,0},X^+_{n,2r+1}\bigr].
\end{align*}
By \eqref{5131}, we obtain
\begin{align*}
&\ad\bigl(X^+_{n-1,0}\bigr)\bigl(-X^+_{n,2r+1}+\bigl[X^+_{n-1,0},X^+_{n,2r+1}\bigr]\bigr)\\
&\qquad{}=-\bigl[X^+_{n-1,0},X^+_{n,2r+1}\bigr]+\bigl[X^+_{n-1,0},\bigl[X^+_{n-1,0},X^+_{n,2r+1}\bigr]\bigr] =-\bigl[X^+_{n-1,0},X^+_{n,2r+1}\bigr]+0,
\\
& \ad\bigl(X^+_{n-1,0}\bigr)\bigl(-[X^+_{n-1,0},X^+_{n,2r+1}]\bigr)=0.
\end{align*}
Thus, we have obtained $\tau_{n-1}\bigl(X^+_{n,2r+1}\bigr)=-X^+_{n,2r+1}$.

Next, we show the $-$ case. By \eqref{5114} and \eqref{5119}, we have
\begin{align*}
\exp\bigl(\ad\bigl(X^+_{n-1,0}\bigr)\bigr)X^-_{n,2r+1}&=X^-_{n,2r+1}+\bigl[X^+_{n-1,0},X^-_{n,2r+1}\bigr]+\frac{1}{2}\bigl[X^+_{n-1,0},\bigl[X^+_{n-1,0},X^-_{n,2r+1}\bigr]\bigr]\\
&=X^-_{n,2r+1}-X^-_{n+1,2r+1}+0.
\end{align*}
By \eqref{5119}, we obtain
\begin{align*}
\bigl(-\ad\bigl(X^-_{n-1,0}\bigr)\bigr)\bigl(X^-_{n,2r+1}-X^-_{n+1,2r+1}\bigr)&=-\bigl[X^-_{n-1,0},X^-_{n,2r+1}\bigr]+\bigl[X^-_{n-1,0},X^-_{n+1,2r+1}\bigr]\\
&=-\bigl[X^-_{n-1,0},X^-_{n,2r+1}\bigr]-2X^-_{n,2r+1}.
\end{align*}
By \eqref{5131}, we have
\begin{align*}
-\frac{1}{2}\ad\bigl(X^-_{n-1,0}\bigr)\bigl(-\bigl[X^-_{n-1,0},X^-_{n,2r+1}\bigr]-2X^-_{n,2r+1}\bigr)=0+\bigl[X^-_{n-1,0},X^-_{n,2r+1}\bigr].
\end{align*}
Then, we obtain
\begin{align*}
\exp\bigl(-\ad\bigl(X^-_{n-1,0}\bigr)\bigr)\exp\bigl(\ad\bigl(X^+_{n-1,0}\bigr)\bigr)X^-_{n,2r+1} =-X^-_{n,2r+1}-X^-_{n+1,2r+1}.
\end{align*}
By \eqref{5114} and \eqref{5119}, we obtain
\begin{align*}
& \ad\bigl(X^+_{n-1,0}\bigr)\bigl(-X^-_{n,2r+1}-X^-_{n+1,2r+1}\bigr) =X^-_{n+1,2r+1}-0,
\\
& \frac{1}{2}\bigl(\ad\bigl(X^+_{n-1,0}\bigr)\bigr)\bigl(X^-_{n+1,2r+1}\bigr)=0.
\end{align*}
Thus, we have obtained $\tau_{n-1}\bigl(X^-_{n,2r+1}\bigr)=-X^-_{n,2r+1}$.
\end{proof}
\begin{Corollary}\label{Cor25}
The following relations hold:{\samepage
\begin{align*}
\tau_i(f_{j,k}^{r,s})&=\begin{cases}
-f_{k,j}^{r,s}&\text{if }(j,k)=(i,i+1),(i+1,i),\\
(-1)^rf_{j,k}^{r,s}&\text{if }(j,k)=(i,-(i+1)),(-(i+1),i),(-(i+1),i),(-i,(i+1)),\\
-f_{j,k}^{r,s}&\text{if }j=i,k\neq\pm(i+1),\\
-f_{-i-1,k}^{r,s}&\text{if }j=-i,k\neq(i+1),\\
f_{i,k}^{r,s}&\text{if }j=i+1,k\neq \pm i,\\
f_{-i,k}^{r,s}&\text{if }j=-(i+1)\neq \pm i,\\
-f_{j,i+1}^{r,s}&\text{if }k=i,j\neq\pm(i+1),\\
-f_{j,-i-1}^{r,s}&\text{if }k=-i,j\neq(i+1),\\
f_{j,i}^{r,s}&\text{if }k=i+1,j\neq i,\\
f_{j,-i}^{r,s}&\text{if }k=-(i+1),j\neq i,\\
f_{j,k}^{r,s}&\text{otherwise},
\end{cases}\\
\tau_n(f_{j,k}^{r,s})&=\begin{cases}
(-1)^rf_{j,k}^{r,s}&\text{if }(j,k)=(n-1,n),(n,n-1),\\
-f_{k,j}^{r,s}&\text{if }(j,k)=(n-1,-n),(n,-(n-1)),(-(n-1),n),\\
-f_{-n,k}^{r,s}&\text{if }j=n-1,k\neq\pm n,\\
-f_{n,k}^{r,s}&\text{if }j=-(n-1),k\neq n,\\
f_{-(n-1),k}&\text{if }j=n,k\neq \pm n-1,\\
f_{(n-1),k}^{r,s}&\text{if }j=-n\neq \pm n-1,\\
-f_{j,-n}^{r,s}&\text{if }k=n-1,j\neq\pm n,\\
-f_{j,n}^{r,s}&\text{if }k=-(n-1),j\neq n,\\
f_{j,-(n-1)}^{r,s}&\text{if }k=n,j\neq n-1,\\
f_{j,(n-1)}^{r,s}&\text{if }k=-n,j\neq n-1,\\
f_{j,k}^{r,s}&\text{otherwise},
\end{cases}\\
\tau_0(f_{j,k}^{r,s})&=\begin{cases}
(-1)^rf_{j,k}^{r,s}&\text{if }(j,k)=(1,2),(j,k)=(2,1),\\
-f_{k,j}^{r-2,s}&\text{if }(j,k)=(1,-2),(1,-2),\\
-f_{k,j}^{r+2,s}&\text{if }(j,k)=(-2,1),(-1,2),\\
f_{-2,j}^{r+2,s}&\text{if }j=1,k\neq\pm 2,\\
f_{2,k}^{r-2,s}&\text{if }j=-1,k\neq 2,\\
-f_{-1,k}^{r+1,s}&\text{if }j=2,k\neq \pm 1,\\
-f_{1,k}^{r-1,s}&\text{if }j=-2\neq 1,\\
f_{j,-2}^{r-1,s}&\text{if }k=1,j\neq\pm 2,\\
f_{j,2}^{r+1,s}&\text{if }k=-1,j\neq 2,\\
-f_{j,-1}^{r-1,s}&\text{if }k=2,j\neq 1,\\
-f_{j,1}^{r+1,s}&\text{if }k=-2,j\neq 1,\\
f_{j,k}^{r,s}&\text{otherwise},
\end{cases}
\end{align*}
for $j\neq \pm k$ and either $j$ or $k$ is positive.}
\end{Corollary}

By Lemma~\ref{Weyl}, we find that $\tau_i\bigl(\mathfrak{D}(2n)_\alpha^k\bigr)$ is contained in $\mathfrak{D}(2n)_{s_i(\alpha)}^k$
In the same way as~\cite[Proposition~3.2]{MRY}, we obtain the following theorem.
\begin{Theorem}\label{Thm128}
We find that
\[\begin{cases}
\dim \mathfrak{D}(2n)_\alpha^k=1&\text{if }\alpha\in\Delta^{{\rm ex}}\setminus\{k\delta\mid k\in\mathbb{Z}\setminus\{0\}\},\\
\dim \mathfrak{D}(2n)_\alpha^k=0&\text{if }\alpha\not\in\Delta^{{\rm ex}}.
\end{cases}
\]
\end{Theorem}
By Theorem~\ref{Thm128}, the kernel of $\pi$ is contained in $\bigoplus_{r\in\mathbb{Z}} \mathfrak{D}(2n)_{r\delta}^k$. Since $\bigl[X^\pm_{i,0},\mathfrak{D}(2n)_{r\delta}^k\bigr]\cap \mathfrak{D}(2n)_{r\delta}^k$ is empty by Lemma~\ref{Lem128}, the kernel of $\pi$ is contained in the center of $D(2n)$.
Thus, we find that $\mathfrak{D}(2n)$ is a central extension of $\mathfrak{sl}(2n)\bigl[u^{\pm1},v\bigr]^\tau$.
\section[Embedding of D(2n) into A(2n)]{Embedding of $\boldsymbol{D(2n)}$ into $\boldsymbol{A(2n)}$}\label{sec:6}
In this section, we show that the associative algebra $\mathfrak{D}(2n)$ is a subalgebra of the associative algebra $A(2n)$.
By a direct computation, we can construct a homomorphism
$\phi\colon \mathfrak{D}(2n)\to A(2n)$
defined by
\begin{gather*}
H_{i,r}=\begin{cases}
f^r_{i,i}v^r-f^r_{i+1,i+1}v^r&\text{if }1\leq i\leq n-1,\\
f^r_{n-1,n-1}v^r+f^r_{n,n}v^r&\text{if }i=n\text{ and $r$ is even},\\
-\bigl(f^r_{1,1}v^r+f^r_{2,2}v^r\bigr)+2u^{-1}v^{r}du&\text{if }i=0\text{ and $r$ is even},
\end{cases}\\
X^+_{i,r}=\begin{cases}
f^r_{i,i+1}v^r&\text{if }1\leq i\leq n-1,\\
f^r_{n-1,-n}v^r&\text{if }i=n,\\
(-1)^rf^r_{-2,1}uv^r&\text{if }i=0,\\
f^r_{n,-n}v^r&\text{if }i=n+1\text{ and $r$ is odd},\\
f^r_{-1,1}uv^r&\text{if }i=-1\text{ and $r$ is odd},
\end{cases}\\
X^-_{i,r}=\begin{cases}
f^r_{i+1,i}v^r&\text{if }1\leq i\leq n-1,\\
f^r_{-n,n-1}v^r&\text{if }i=n,\\
(-1)^rf^r_{1,-2}u^{-1}v^r&\text{if }i=0,\\
f^r_{-n,n}v^r&\text{if }i=n+1\text{ and $r$ is odd},\\
f^r_{1,-1}u^{-1}v^r&\text{if }i=-1\text{ and $r$ is odd}.
\end{cases}
\end{gather*}

\begin{Theorem}\label{inj}
The homomorphism $\phi$ is injective.
\end{Theorem}

This section is devoted to the proof of Theorem~\ref{inj}.
By Theorem~\ref{Thm128}, we can define $\smash{\widetilde{f}_{i,j}^{r,s}\in \mathfrak{D}(2n)}$ for $i\neq j$ in the same way as $f_{i,j}^{r,s}\in\mathfrak{sl}(2n)\bigl[u^{\pm1},v\bigr]^\tau$. By Theorem~\ref{Thm128}, we obtain the relations
\begin{gather*}
\bigl[H_{i,x},\widetilde{f}_{p,q}^{r,s}\bigr]\nonumber\\
=\begin{cases}
(\delta_{i,p}-\delta_{i+1,p}-\delta_{i,q}+\delta_{i+1,q}-(-1)^x(\delta_{-i,p}-\delta_{-i,q}-\delta_{-i-1,p}+\delta_{-i-1,q}))\widetilde{f}_{p,q}^{r,s+x}\\
\quad \text{if }1\leq i\leq n-1,\\
(\delta_{n-1,p}+\delta_{n,p}-\delta_{n-1,q}-\delta_{n,q}-(-1)^x(\delta_{-(n-1),p}-\delta_{-(n-1),q}+\delta_{-n,p}-\delta_{-n,q}))\widetilde{f}_{p,q}^{r,s+x}\\
\quad \text{if }i=n,\\
-(\delta_{1,p}+\delta_{2,p}-\delta_{1,q}-\delta_{2,q}-(-1)^x(\delta_{-1,p}-\delta_{-1,q}+\delta_{-2,p}-\delta_{-2,q}))\widetilde{f}_{p,q}^{r,s+x}\\
\quad \text{if }i=0,
\end{cases}\\
\bigl[\widetilde{f}_{i,j}^{r_1,s_1},\widetilde{f}_{p,q}^{r_2,s_2}\bigr]
=\delta_{j,p}\widetilde{f}_{i,q}^{r_1+r_2,s_1+s_2}-\delta_{i,q}\widetilde{f}_{p,j}^{r_1+r_2,s_1+s_2}\nonumber\\
\quad{}
-(-1)^{s_1}\delta_{i,-p}\widetilde{f}_{-j,q}^{r_1+r_2,s_1+s_2}+(-1)^{s_1}\delta_{j,-q}\widetilde{f}_{p,-i}^{r_1+r_2,s_1+s_2}\qquad \text{if }(i,j)\neq(q,p),(-p,-q).
\end{gather*}

By Theorem~\ref{Thm128}, we find that the kernel of $\Phi$ is contained in $\bigoplus_{r\neq0,k\in\mathbb{Z}}\mathfrak{D}(2n)_{r\delta}^k$.
\begin{Lemma}
For $r\geq0$, the subspace $\mathfrak{D}(2n)_{r\delta}^k$ is spanned by
\begin{align*}
a(r,k)={}&\bigl\{\bigl[X^+_{i,0},\widetilde{f}_{i+1,i}^{r,k}\bigr]\mid1\leq i\leq n-1\bigr\}\cup\bigl\{\bigl[X^+_{n,0},\widetilde{f}_{-n,n-1}^{r,k}\bigr]\mid1\leq i\leq n-1\bigr\}\\
&{\cup}\,\bigl\{\bigl[X^+_{0,s},\widetilde{f}_{1,-2}^{r,k}\bigr]\mid s=0,1\bigr\}\cup\bigl\{\bigl[X^+_{1,1},\widetilde{f}_{2,1}^{r,k-1}\bigr]-\bigl[X^+_{1,0},\widetilde{f}_{2,1}^{r,k}\bigr]\bigr\}.
\end{align*}
\end{Lemma}
\begin{proof}
By Theorem~\ref{Thm128}, we find that $\mathfrak{D}(2n)_{r\delta}^k$ is spanned by the set
\begin{align*}
&\bigl\{\bigl[X^+_{i,s},\widetilde{f}_{i+1,i}^{r,v}\bigr]\mid1\leq i\leq n-1,s+v=k\bigr\}\cup\bigl\{\bigl[X^+_{n,s},\widetilde{f}_{-n,n-1}^{r,v}\bigr]\mid s+v=k\bigr\}\\
&\qquad{}\cup\bigl\{\bigl[X^+_{0,s},\widetilde{f}_{1,-2}^{r-1,v}\bigr]\mid s+v=k\bigr\}.
\end{align*}
Since $\mathfrak{D}(2n)$ is a central extension of $\mathfrak{sl}(2n)\bigl[u^{\pm1},v\bigr]^\tau$, we find that
\begin{gather*}
\bigl[X^+_{i,s},\widetilde{f}_{i+1,i}^{r,v}\bigr]-\bigl[X^+_{i,0},\widetilde{f}_{i+1,i}^{r,s+v}\bigr],\\
\bigl[X^+_{n,s},\widetilde{f}_{-n,n-1}^{r,v}\bigr]-\bigl[X^+_{n,0},\widetilde{f}_{-n,n-1}^{r,s+v}\bigr],\\
\bigl[X^+_{0,s},\widetilde{f}_{1,-2}^{r-1,v}\bigr]-\bigl[X^+_{0,0},\widetilde{f}_{1,-2}^{r-1,s+v}\bigr]
\end{gather*}
are central elements of $\mathfrak{D}(2n)$. Thus, we obtain
\begin{align}
0&=\bigl[H_{i,1},\bigl[X^+_{i,s},\widetilde{f}_{i+1,i}^{r,v}\bigr]\bigr]-\bigl[H_{i,1},\bigl[X^+_{i,0}, \widetilde{f}_{i+1,i}^{r,s+v}\bigr]\bigr]=2\bigl[X^+_{i,s+1},\widetilde{f}_{i+1,i}^{r,v}\bigr] -2\bigl[X^+_{i,s},\widetilde{f}_{i+1,i}^{r,v+1}\bigr]\nonumber\\
&\quad-2\bigl[X^+_{i,1},\widetilde{f}_{i+1,i}^{r,s+v}\bigr]+2\bigl[X^+_{i,0},\widetilde{f}_{i+1,i}^{r,s+v+1}\bigr]\qquad \text{for }1\leq i\leq n-1,\label{eq92}\\
0&=\bigl[H_{n-1,1},\bigl[X^+_{n,s},\widetilde{f}_{-n,n-1}^{r,v}\bigr]\bigr]-\bigl[H_{n-1,1}, \bigl[X^+_{n,0},\widetilde{f}_{-n,n-1}^{r,s+v}\bigr]\bigr]\nonumber\\%\label{eq92-1}
&=2\bigl[X^+_{n,s+1},\widetilde{f}_{-n,n-1}^{r,v}\bigr]-2\bigl[X^+_{n,s},\widetilde{f}_{-n,n-1}^{r,v+1}\bigr] -2\bigl[X^+_{n,1},\widetilde{f}_{-n,n-1}^{r,s+v}\bigr]+2\bigl[X^+_{n,0},\widetilde{f}_{-n,n-1}^{r,s+v+1}\bigr],\nonumber\\
0&=\bigl[H_{1,1},\bigl[X^+_{0,s},\widetilde{f}_{1,-2}^{r-1,v}\bigr]\bigr]-\bigl[H_{1,1},\bigl[X^+_{0,0},\widetilde{f}_{1,-2}^{r-1,s+v}\bigr]\bigr]\nonumber\\%\label{eq93}
&=2\bigl[X^+_{0,s+1},\widetilde{f}_{1,-2}^{r-1,v}\bigr]-2\bigl[X^+_{0,s},\widetilde{f}_{1,-2}^{r-1,v+1}\bigr] -2\bigl[X^+_{0,1},\widetilde{f}_{1,-2}^{r-1,s+v}\bigr]+2\bigl[X^+_{0,0},\widetilde{f}_{1,-2}^{r-1,s+v+1}\bigr].\nonumber
\end{align}
Thus, $\mathfrak{A}(2n)_{r\delta}^k$ is spanned by
\begin{align*}
&\{\bigl[X^+_{i,s},\widetilde{f}_{i+1,i}^{r,k-s}\bigr]\mid1\leq i\leq n-1,\,s=0,1\}\\
& \qquad\cup\{\bigl[X^+_{n,s},\widetilde{f}_{-n,n-1}^{r,k-s}\bigr]\mid1\leq i\leq n-1,s=0,1\}\cup\{\bigl[X^+_{0,s},\widetilde{f}_{1,-2}^{r-1,k-s}\bigr]\mid s=0,1\}.
\end{align*}
By Corollary~\ref{Cor25} and a direct computation, we find that, by composing $\tau_i$, we can construct~$\widetilde{\tau}_1$,~$\widetilde{\tau}_2$ satisfying that
\begin{alignat*}{3}
& \widetilde{\tau}_1\bigl(X^+_{i,k}\bigr)=X^+_{1,k},\qquad && \widetilde{\tau}_1\bigl(\widetilde{f}_{i,i+1}^{r,k}\bigr)=\widetilde{f}_{2,1}^{r,k},&\\
&\widetilde{\tau}_2\bigl(X^+_{0,k}\bigr)=X^+_{1,k},\qquad && \widetilde{\tau}_2\bigl(\widetilde{f}_{1,-2}^{r-1,k}\bigr)=\widetilde{f}_{2,1}^{r,k},&\\
&\widetilde{\tau}_3\bigl(X^+_{n,k}\bigr)=X^+_{1,k},\qquad && \widetilde{\tau}_3\bigl(\widetilde{f}_{-n,n-1}^{r,k}\bigr)=\widetilde{f}_{2,1}^{r,k}&
\end{alignat*}
for $1\leq i\leq n$. Then, we obtain
\begin{gather}
\bigl[X^+_{1,1},\widetilde{f}_{2,1}^{r,k-1}\bigr]-\bigl[X^+_{1,0},\widetilde{f}_{2,1}^{r,k}\bigr]
=\bigl[\widetilde{\tau}_1\bigl(X^+_{i,1}\bigr),\widetilde{\tau}_1\bigl(\widetilde{f}_{i+1,i}^{r,k-1}\bigr)\bigr]-\bigl[\widetilde{\tau}_1\bigl(X^+_{i,0}\bigr),\widetilde{\tau}_1\bigl(\widetilde{f}_{i+1,i}^{r,k}\bigr)\bigr]\nonumber\\
\hphantom{\bigl[X^+_{1,1},\widetilde{f}_{2,1}^{r,k-1}\bigr]-\bigl[X^+_{1,0},\widetilde{f}_{2,1}^{r,k}\bigr]}{}
=\bigl[X^+_{i,1},\widetilde{f}_{i+1,i}^{r,k-1}\bigr]-\bigl[X^+_{i,0},\widetilde{f}_{i+1,i}^{r,k}\bigr],\label{eq94}\\
\bigl[X^+_{1,1},\widetilde{f}_{2,1}^{r,k-1}\bigr]-\bigl[X^+_{1,0},\widetilde{f}_{2,1}^{r,k}\bigr]
=\bigl[\widetilde{\tau}_2\bigl(X^+_{0,1}\bigr),\widetilde{\tau}_2\bigl(\widetilde{f}_{1,-2}^{r-1,k-1}\bigr)\bigr]-\bigl[\widetilde{\tau}_2\bigl(X^+_{0,0}\bigr),\widetilde{\tau}_2\bigl(\widetilde{f}_{1,-2}^{r-1,k}\bigr)\bigr]\nonumber\\
\hphantom{\bigl[X^+_{1,1},\widetilde{f}_{2,1}^{r,k-1}\bigr]-\bigl[X^+_{1,0},\widetilde{f}_{2,1}^{r,k}\bigr]}{}
=\bigl[X^+_{0,1},\widetilde{f}_{1,-2}^{r-1,k-1}\bigr]-\bigl[X^+_{0,0},\widetilde{f}_{1,-2}^{r-1,k}\bigr],\nonumber\\%\label{eq95}
\bigl[X^+_{1,1},\widetilde{f}_{2,1}^{r,k-1}\bigr]-\bigl[X^+_{1,0},\widetilde{f}_{2,1}^{r,k}\bigr]
=\bigl[\widetilde{\tau}_3\bigl(X^+_{n,1}\bigr),\widetilde{\tau}_3\bigl(\widetilde{f}_{-n,n-1}^{r,k-1}\bigr)\bigr]
-\bigl[\widetilde{\tau}_3\bigl(X^+_{n,0}\bigr),\widetilde{\tau}_3\bigl(\widetilde{f}_{-n,n-1}^{r,k}\bigr)\bigr]\nonumber\\
\hphantom{\bigl[X^+_{1,1},\widetilde{f}_{2,1}^{r,k-1}\bigr)\bigr]-\bigl[X^+_{1,0},\widetilde{f}_{2,1}^{r,k}\bigr)\bigr]}{}
=\bigl[X^+_{n,1},\widetilde{f}_{-n,n-1}^{r,k-1}\bigr]-\bigl[X^+_{n,0},\widetilde{f}_{-n,n-1}^{r,k}\bigr],\nonumber%\label{eq95-1}
\end{gather}
because the right-hand sides are central elements of $D(2n)$.
\end{proof}

Now, we begin the proof of Theorem~\ref{inj}. We only consider the case that $r>0$. When $k$ is even, the images of $a(r,k)$ via $\Phi$ are linearly independent. However, in the case that $k$ is odd, we obtain
\begin{equation*}
\bigl[\Phi\bigl(X^+_{1,1}\bigr),\Phi\bigl(\widetilde{f}_{2,1}^{r,k}\bigr)\bigr]-\bigl[\Phi\bigl(X^+_{1,0}\bigr),\Phi\bigl(\widetilde{f}_{2,1}^{r,k}\bigr)\bigr]=0.
\end{equation*}
Thus, we need to show \smash{$\bigl[X^+_{1,1},\widetilde{f}_{2,1}^{r,2s}\bigr]-\bigl[X^+_{1,0},\widetilde{f}_{2,1}^{r,2s+1}\bigr]=0$}.
The case $s=0$ is precisely the defining relation.
Next, we consider the case $s\geq1$. By a direct computation, we obtain
\begin{align*}
\bigl[X^+_{n,2},f_{-n,n-1}^{u,2s-1}\bigr]&=-\frac{1}{2}\bigl[X^+_{n,2},\bigl[X^+_{n-1,0},f_{-(n-1),n-1}^{u,2s-1}\bigr]\bigr]\\
&=-\frac{1}{2}\bigl[X^+_{n-1,0},\bigl[X^+_{n,2},f_{-(n-1),n-1}^{u,2s-1}\bigr]\bigr]=\bigl[X^+_{n-1,0},f_{n,n-1}^{u,2s+1}\bigr]
\end{align*}
and
\begin{align*}
\bigl[X^+_{n,0},f_{-n,n-1}^{u,2s+1}\bigr]&=-\frac{1}{2}\bigl[X^+_{n,0},\bigl[X^+_{n-1,2},f_{-(n-1),n-1}^{u,2s-1}\bigr]\bigr]\\
&=-\frac{1}{2}\bigl[X^+_{n-1,2},\bigl[X^+_{n,0},f_{-(n-1),n-1}^{u,2s-1}\bigr]\bigr]=\bigl[X^+_{n-1,2},f_{n,n-1}^{u,2r-1}\bigr].
\end{align*}
Thus, we obtain
\begin{equation}
\bigl[X^+_{n,2},f_{-n,n-1}^{u,2s-1}\bigr]-\bigl[X^+_{n,0},f_{-n,n-1}^{u,2s+1}\bigr]=\bigl[X^+_{n-1,0},f_{n,n-1}^{u,2s+1}\bigr]-\bigl[X^+_{n-1,2},f_{n,n-1}^{u,2s-1}\bigr].\label{eq96}
\end{equation}
By \eqref{eq92}, \eqref{eq96} and \eqref{eq94}, we find the relations
\begin{gather*}
\bigl[X^+_{n,2},f_{-n,n-1}^{u,2s-1}\bigr]-\bigl[X^+_{n,0},f_{-n,n-1}^{u,2s+1}\bigr]=0,\\
\bigl[X^+_{1,1},\widetilde{f}_{2,1}^{2r,2s-1}\bigr]-\bigl[X^+_{1,0},\widetilde{f}_{2,1}^{2r+1,k}\bigr]=0.
\end{gather*}
We complete the proof of Theorem~\ref{inj}.

\begin{thebibliography}{99}
\footnotesize\itemsep=0pt

\bibitem{BR1}
Belliard S., Regelskis V., Drinfeld~{J} presentation of twisted {Y}angians,
 \href{https://doi.org/10.3842/SIGMA.2017.011}{\textit{SIGMA}} \textbf{13}
 (2017), 011, 35~pages,
 \href{https://arxiv.org/abs/1401.2143}{arXiv:1401.2143}.

\bibitem{Bro}
Brown J., Twisted {Y}angians and finite {$W$}-algebras,
 \href{https://doi.org/10.1007/s00031-008-9041-x}{\textit{Transform. Groups}}
 \textbf{14} (2009), 87--114,
 \href{https://arxiv.org/abs/0710.2918}{arXiv:0710.2918}.

\bibitem{BK}
Brundan J., Kleshchev A., Shifted {Y}angians and finite {$W$}-algebras,
 \href{https://doi.org/10.1016/j.aim.2004.11.004}{\textit{Adv. Math.}}
 \textbf{200} (2006), 136--195,
 \href{https://arxiv.org/abs/math/0407012}{arXiv:math/0407012}.

\bibitem{DKV}
De~Sole A., Kac V.G., Valeri D., A~{L}ax type operator for quantum finite
 {$W$}-algebras,
 \href{https://doi.org/10.1007/s00029-018-0439-6}{\textit{Selecta
 Math.~(N.S.)}} \textbf{24} (2018), 4617--4657,
 \href{https://arxiv.org/abs/1707.03669}{arXiv:1707.03669}.

\bibitem{D1}
Drinfeld V.G., Hopf algebras and the quantum {Y}ang--{B}axter equation,
 \textit{Soviet Math. Dokl.} \textbf{32} (1985), 256--258.

\bibitem{D2}
Drinfeld V.G., A~new realization of {Y}angians and of quantum affine algebras,
 \textit{Soviet Math. Dokl.} \textbf{36} (1988), 212--216.

\bibitem{Gu1}
Guay N., Affine {Y}angians and deformed double current algebras in type~{$A$},
 \href{https://doi.org/10.1016/j.aim.2006.08.007}{\textit{Adv. Math.}}
 \textbf{211} (2007), 436--484.

\bibitem{GNW}
Guay N., Nakajima H., Wendlandt C., Coproduct for {Y}angians of affine
 {K}ac--{M}oody algebras,
 \href{https://doi.org/10.1016/j.aim.2018.09.013}{\textit{Adv. Math.}}
 \textbf{338} (2018), 865--911,
 \href{https://arxiv.org/abs/1701.05288}{arXiv:1701.05288}.

\bibitem{HU}
Harako S., Ueda M., A~new presentation of the twisted {Y}angian of type~{$D$},
 \href{https://arxiv.org/abs/2507.00350}{arXiv:2507.00350}.

\bibitem{K}
Kac V.G., Infinite-dimensional {L}ie algebras, 3rd ed.,
 \href{https://doi.org/10.1017/CBO9780511626234}{Cambridge University Press},
 Cambridge, 1990.

\bibitem{Kassel}
Kassel C., K\"ahler differentials and coverings of complex simple {L}ie
 algebras extended over a commutative algebra,
 \href{https://doi.org/10.1016/0022-4049(84)90040-9}{\textit{J.~Pure Appl.
 Algebra}} \textbf{34} (1984), 265--275.

\bibitem{L}
Lu K., Minimalistic presentation and coideal structure of twisted {Y}angians,
 \href{https://doi.org/10.1007/s00220-026-05615-3}{\textit{Comm. Math. Phys.}}
 \textbf{407} (2026), 99, 38~pages,
 \href{https://arxiv.org/abs/2511.07136}{arXiv:2511.07136}.

\bibitem{LPTTW}
Lu K., Peng Y.N., Tappeiner L., Topley L., Wang W., Shifted twisted {Y}angians
 and finite {$W$}-algebras of classical type,
 \href{https://arxiv.org/abs/2505.03316}{arXiv:2505.03316}.

\bibitem{LWZ0}
Lu K., Wang W., Zhang W., A~{D}rinfeld type presentation of twisted {Y}angians,
 \href{https://doi.org/10.1090/ert/702}{\textit{Represent. Theory}}
 \textbf{29} (2025), 838--870,
 \href{https://arxiv.org/abs/2308.12254}{arXiv:2308.12254}.

\bibitem{LWZ}
Lu K., Wang W., Zhang W., Affine {$\imath$}quantum groups and twisted
 {Y}angians in {D}rinfeld presentations,
 \href{https://doi.org/10.1007/s00220-025-05263-z}{\textit{Comm. Math. Phys.}}
 \textbf{406} (2025), 98, 36~pages,
 \href{https://arxiv.org/abs/2406.05067}{arXiv:2406.05067}.

\bibitem{LWZ1}
Lu K., Zhang W., A~{D}rinfeld type presentation of twisted {Y}angians of
 quasi-split type,
 \href{https://doi.org/10.1142/S021919972650001X}{\textit{Commun. Contemp.
 Math.}}, {t}o appear,
 \href{https://arxiv.org/abs/2408.06981}{arXiv:2408.06981}.

\bibitem{MNT}
Matsuo A., Nagatomo K., Tsuchiya A., Quasi-finite algebras graded by
 {H}amiltonian and vertex operator algebras, in Moonshine: the First Quarter
 Century and Beyond, \textit{London Math. Soc. Lecture Note Ser.}, Vol.~372,
 \href{https://doi.org/10.1017/CBO9780511730054.015}{Cambridge University
 Press}, Cambridge, 2010, 282--329,
 \href{https://arxiv.org/abs/math/0505071}{arXiv:math/0505071}.

\bibitem{MRY}
Moody R.V., Rao S.E., Yokonuma T., Toroidal {L}ie algebras and vertex
 representations, \href{https://doi.org/10.1007/BF00147350}{\textit{Geom.
 Dedicata}} \textbf{35} (1990), 283--307.

\bibitem{O}
Olshanskii G.I., Twisted {Y}angians and infinite-dimensional classical {L}ie
 algebras, in Quantum Groups ({L}eningrad, 1990), \textit{Lecture Notes in
 Math.}, Vol.~1510, \href{https://doi.org/10.1007/BFb0101183}{Springer},
 Berlin, 1992, 104--119.

\bibitem{U77}
Ueda M., Twisted affine {Y}angian and rectangular {$W$}-algebra of type~{$D$},
 \href{https://doi.org/10.1007/s10468-025-10348-3}{\textit{Algebr. Represent.
 Theory}} \textbf{28} (2025), 1195--1224,
 \href{https://arxiv.org/abs/2107.09999}{arXiv:2107.09999}.

\bibitem{U78}
Ueda M., Suggestion of a~definition of the twisted affine Yangian of
 type~{$D$}, \href{https://arxiv.org/abs/2512.06340}{arXiv:2512.06340}.

\end{thebibliography} 
\end{document}
